\documentclass[tikz,border=8pt]{standalone}
\usepackage{tikz}
\usepackage{amsmath}
\usepackage{amssymb}
\usepackage{pgfplots}
\pgfplotsset{compat=1.18}
\usepackage{ctex}
\usetikzlibrary{calc, positioning, arrows.meta, matrix}

\begin{document}
\begin{tikzpicture}[
  box/.style={
    draw=gray!60, thin, rounded corners=4pt,
    fill=white, inner sep=10pt,
    font=\small
  },
  arr/.style={->, >=Stealth, thick, gray!70!black},
  steplabel/.style={font=\scriptsize, fill=white, draw=gray!50, thin,
                    rounded corners=2pt, inner sep=3pt, align=center}
]

%% ===== 框1：原增广矩阵 =====
\node[box] (mat1) at (0,0) {%
  \begin{tabular}{c}
  \textbf{原增广矩阵}\\[4pt]
  $\left[\begin{array}{ccc|c}
    1 & 2 & 1 & 5 \\
    2 & 5 & 3 & 11 \\
    1 & 1 & 0 & 3
  \end{array}\right]$
  \end{tabular}
};

%% ===== 框2：上三角形式 =====
\node[box] (mat2) at (0,-4.8) {%
  \begin{tabular}{c}
  \textbf{上三角矩阵（行阶梯形）}\\[4pt]
  $\left[\begin{array}{ccc|c}
    1 & 2 & 1 & 5 \\
    0 & 1 & 1 & 1 \\
    0 & 0 & 1 & 2
  \end{array}\right]$
  \end{tabular}
};

%% ===== 框3：RREF =====
\node[box] (mat3) at (0,-9.6) {%
  \begin{tabular}{c}
  \textbf{最简行阶梯形（RREF）}\\[4pt]
  $\left[\begin{array}{ccc|c}
    1 & 0 & 0 & 1 \\
    0 & 1 & 0 & -1 \\
    0 & 0 & 1 & 2
  \end{array}\right]$
  \end{tabular}
};

%% ===== 箭头1: mat1 -> mat2 =====
\draw[arr] (mat1.south) -- node[steplabel, right=8pt] {
  $R_2 \leftarrow R_2 - 2R_1$ \\[2pt]
  $R_3 \leftarrow R_3 - R_1$
} (mat2.north);

%% ===== 箭头2: mat2 -> mat3 =====
\draw[arr] (mat2.south) -- node[steplabel, right=8pt] {
  $R_1 \leftarrow R_1 - R_3$ \\[2pt]
  $R_1 \leftarrow R_1 - 2R_2$ \\[2pt]
  $R_2 \leftarrow R_2 - R_3$
} (mat3.north);

%% ===== 底部注释：主元与自由变量 =====
\node[font=\small, fill=white, draw=gray!50, thin, rounded corners=2pt,
      inner sep=6pt, align=center] at (0,-11.4) {
  \textbf{主元（Pivot）}：$x_1,\,x_2,\,x_3$（共 3 个）\quad
  \textbf{自由变量}：无 \quad
  \textbf{唯一解}：$x_1=1,\;x_2=-1,\;x_3=2$
};

%% 整体标题
\node[font=\large\bfseries] at (0, 1.6) {高斯消元法过程示意};

\end{tikzpicture}
\end{document}
